\documentclass[11pt]{article}
\usepackage[paperwidth=8.5in,paperheight=11in,margin=1in]{geometry}
\usepackage{xr}
\usepackage{amsmath}
\externaldocument{paper}

% For section headers starting with S
\renewcommand{\thesection}{S.\arabic{section}}
\renewcommand{\thesubsection}{\thesection.\arabic{subsection}}

% Equation numbering in supplementary style
\makeatletter
\def\tagform@#1{\maketag@@@{(S\ignorespaces#1\unskip\@@italiccorr)}}
\makeatother

\newcommand{\myfigurename}{Figure}
\newcommand{\myref}[1]{\ref{#1}}
\newcommand{\suppref}[1]{S\ref{#1}}
\usepackage{amsmath,amssymb}
\DeclareMathOperator{\EX}{\mathrm{E}}
\DeclareMathOperator{\Var}{\mathrm{Var}}
\DeclareMathOperator{\Cov}{\mathrm{Cov}}
\usepackage{graphicx}

\title{Supporting Information}
\author{}

\begin{document}
\maketitle

\section{Methods}

\subsection{The EAGGL model}

EAGGL (Enrichment Analysis Gene Grouping and LLMs) is the latent-structure stage that follows PIGEAN. Whereas PIGEAN estimates the probability that each gene is relevant to a trait and the degree to which each gene annotation predicts trait-relevance, the goal of EAGGL is to decompose those trait-associated gene and annotation patterns into a small number of latent mechanisms.

Let genes be indexed by $i = 1 \ldots N$, gene annotations or gene sets by $j = 1 \ldots M$, phenotypes by $t = 1 \ldots T$, and latent mechanisms by $k = 1 \ldots K$. Let $X_{ij} \ge 0$ denote the value of annotation $j$ for gene $i$. In the simplest setting, $X_{ij}$ is binary gene-set membership, although weighted annotations are also permitted. Let $p_i^{(u)} \in [0,1]$ denote the gene-level support for gene $i$ for anchor user $u$, and let $q_j^{(u)} \in [0,1]$ denote the corresponding annotation-level or gene-set-level support for annotation $j$ for anchor user $u$. The index $u = 1 \ldots U$ is a generic anchor index; depending on the workflow, it may refer to the default phenotype, a supplied positive-control pseudo-phenotype, one or more anchor phenotypes, one or more anchor genes, or an anchor gene set.

The basic biological assumption is that trait-relevant genes and trait-relevant annotations co-occur because they load on a small number of shared mechanisms. In the gene-by-annotation workflow family, we therefore posit nonnegative gene weights $W_{ik}$ and nonnegative annotation weights $H_{jk}$ such that
\begin{equation}
X_{ij} \approx \sum_{k=1}^K W_{ik} H_{jk} + \epsilon_{ij},
\label{eq:eaggl_basic_factor}
\end{equation}
where $\epsilon_{ij}$ is an error term. A factor $k$ is interpreted as a latent mechanism when a coherent set of genes has large $W_{ik}$ values and a coherent set of annotations has large $H_{jk}$ values.

The entries of $X$ are not equally informative for a given trait. For example, an annotation that is itself strongly implicated by PIGEAN should contribute more strongly to inference than an annotation with no trait support; likewise, genes with high trait support should contribute more than genes with no support. EAGGL therefore fits the factorization under a weighted Gaussian likelihood:
\begin{equation}
\Pr(X \mid W, H, \Omega)
\propto
\exp\left(
-\frac{1}{2}
\sum_{i=1}^N \sum_{j=1}^M
\Omega_{ij}
\left(
X_{ij} - \sum_{k=1}^K W_{ik} H_{jk}
\right)^2
\right),
\label{eq:eaggl_weighted_likelihood}
\end{equation}
where $\Omega_{ij} \ge 0$ is a cell-specific weight. In the simplest single-trait setting, $\Omega_{ij}=p_iq_j$, where $p_i$ is gene support and $q_j$ is annotation support. With multiple anchor traits, the default \texttt{multi} aggregation uses same-trait additive support,
\begin{equation}
\Omega_{ij}^{\mathrm{multi}}
=
\sum_{t=1}^T p_i^{(t)} q_j^{(t)}.
\label{eq:eaggl_weight_matrix}
\end{equation}
The explicit \texttt{any} aggregation instead uses cell-level noisy-OR over same-trait co-support,
\begin{equation}
\Omega_{ij}^{\mathrm{any}}
=
1 - \prod_{t=1}^T \left(1 - p_i^{(t)}q_j^{(t)}\right).
\label{eq:eaggl_any_weight_matrix}
\end{equation}
Both definitions reduce exactly to $p_iq_j$ with one anchor trait, and a trait with zero support everywhere does not change either weight matrix. The noisy-OR is applied to $p_i^{(t)}q_j^{(t)}$ directly, not to separately aggregated gene and annotation probabilities, so it does not create cross-trait support.

Although EAGGL is often motivated heuristically by factorizing a matrix of the form $X_{ij} p_i q_j$, the weighted likelihood in Eq.~\eqref{eq:eaggl_weighted_likelihood} is the more useful mathematical form. It makes explicit that EAGGL seeks a factorization of the underlying annotation matrix while emphasizing entries most relevant to the trait or anchor set. This also avoids repeatedly materializing a different dense weighted matrix for every anchor user.

\subsection{Fitting the model}

Because the factor matrices are latent, we infer them from the observed nonnegative matrix and the support-derived weights. The factorization is regularized by automatic relevance determination (ARD), which allows the effective number of factors to be learned from the data rather than fixed in advance.

Specifically, for each factor $k$ we introduce a nonnegative relevance scale $\lambda_k$ and use the prior
\begin{equation}
W_{ik} \mid \lambda_k \sim \text{HalfNormal}\left(0, \lambda_k / \phi \right),
\qquad
H_{jk} \mid \lambda_k \sim \text{HalfNormal}\left(0, \lambda_k / \phi \right),
\label{eq:eaggl_halfnormal_prior}
\end{equation}
together with a factor-specific ARD prior
\begin{equation}
\lambda_k \sim \text{InvGamma}(a_0, b_0).
\label{eq:eaggl_ard_prior}
\end{equation}
Here $\phi$ controls the baseline sparsity of the factor loadings, while $a_0$ and $b_0$ control the tendency of entire factors to be shrunk toward low relevance. Small values of $\lambda_k$ correspond to factors that explain little of the observed matrix after accounting for the regularization. ARD is a soft relevance shrinkage mechanism rather than a guarantee that every low-mass column is driven exactly to zero.

Up to additive constants, maximizing the posterior is equivalent to minimizing
\begin{equation}
\mathcal{L}(W,H,\lambda)
=
\frac{1}{2}
\sum_{i=1}^N \sum_{j=1}^M
\Omega_{ij}
\left(
X_{ij} - \sum_{k=1}^K W_{ik} H_{jk}
\right)^2
+
\frac{\phi}{2}
\sum_{k=1}^K
\frac{
\|W_{\cdot k}\|_2^2 + \|H_{\cdot k}\|_2^2 + 2b_0
}{
\lambda_k
}
+
C
\sum_{k=1}^K
\log \lambda_k,
\label{eq:eaggl_objective}
\end{equation}
where
\begin{equation}
C = \frac{N + M}{2} + a_0 + 1.
\label{eq:eaggl_C}
\end{equation}
Thus, EAGGL uses the same basic intuition as other Bayesian nonnegative matrix factorization procedures with ARD: the reconstruction term encourages faithful explanation of the observed matrix, while the ARD penalty favors retaining only factors whose contribution is large enough to overcome whole-factor shrinkage.

The number of retained mechanisms is therefore not simply the user-specified maximum candidate rank. Rather, EAGGL begins from a candidate rank $K$ and then uses both the posterior relevance scales and post-fit loading mass to distinguish interpretable factors from the low-mass ARD tail. In this way, the effective number of mechanisms is learned adaptively while retaining explicit diagnostics for the fitted tail.

\subsection{Single-trait gene-by-annotation factorization}

The default EAGGL workflow operates on the gene-by-annotation matrix $X$ together with a single trait-specific gene support vector $p_i$ and annotation support vector $q_j$. In this case $U=1$ and the weight matrix simplifies to
\begin{equation}
\Omega_{ij} = p_i q_j.
\label{eq:eaggl_single_trait_weights}
\end{equation}
The fitted factor matrices $W$ and $H$ may be interpreted as:
\begin{itemize}
\item $W_{ik}$: the degree to which gene $i$ participates in mechanism $k$;
\item $H_{jk}$: the degree to which annotation $j$ captures mechanism $k$.
\end{itemize}
This is the primary EAGGL model and corresponds to the scientific setting in which one seeks to decompose the trait-associated gene-by-annotation landscape into a small number of mechanisms.

A useful special case is when the input support is not obtained from a standard PIGEAN trait run, but instead from a manually supplied gene list. In that standalone workflow, EAGGL first uses the loaded gene universe from $X$ as the enrichment background and computes a hypergeometric enrichment $p$-value for every annotation or gene set. If $M$ is the number of genes in the loaded universe, $n$ is the number of supplied genes present in that universe, $N_j$ is the size of annotation $j$, and $k_j$ is the overlap between annotation $j$ and the supplied gene list, then
\begin{equation}
p_j^{(\mathrm{hyper})}
=
\Pr\!\left(
\mathrm{Hypergeom}(M, n, N_j) \ge k_j
\right).
\label{eq:eaggl_gene_list_hypergeom}
\end{equation}
After Benjamini--Hochberg correction, annotations with $q_j \le q_{\max}$ are retained. Their annotation-level support is then defined by
\begin{equation}
q_j^{(\mathrm{standalone})}
=
\frac{-\log p_j^{(\mathrm{hyper})}}{\sqrt{N_j}},
\label{eq:eaggl_gene_list_weights}
\end{equation}
while the gene-level support vector is binary:
\[
p_i^{(\mathrm{standalone})}
=
\begin{cases}
1 & \text{if gene } i \text{ is in the supplied list,} \\
0 & \text{otherwise.}
\end{cases}
\]
All genes appearing in any retained annotation are then kept in the final factoring matrix. The factorization itself is unchanged; only the way the gene-level and annotation-level support vectors are constructed differs from the standard PIGEAN-driven path.

\subsection{Gene-by-gene symmetric mechanism discovery}
\label{sec:eaggl_gene_gene_symnmf}

The standard gene-by-annotation EAGGL model discovers mechanisms by factorizing a
weighted gene-by-annotation matrix. This makes annotation rows active observations in
mechanism discovery. As an alternative, EAGGL can discover mechanisms directly in gene
space by first constructing a nonnegative gene-by-gene matrix of pairwise
mechanism-sharing evidence, and then applying a symmetric nonnegative factorization to
that matrix. In this formulation, the latent mechanisms are functions from genes to
probabilities of membership, while annotations are used to construct pairwise evidence
and are projected onto the learned mechanisms afterward.

\paragraph{Latent mechanisms.}
Let $z_{ik}\in\{0,1\}$ indicate whether gene $i$ belongs to latent mechanism $k$.
The fitted loading
\[
W_{ik}\in[0,1]
\]
is interpreted as an approximate posterior probability or membership score for the
event $z_{ik}=1$. We assume that mechanisms are sparse, nonnegative, and approximately
disjoint at the level of a single gene. Under exact disjointness,
\begin{equation}
\Pr\{\text{gene } i \text{ belongs to any discovered mechanism}\}
= \sum_k W_{ik}.
\label{eq:eaggl_gene_gene_gene_union}
\end{equation}
For two genes, the probability that they share at least one mechanism is
\begin{equation}
\Pr\{i,j\text{ share a discovered mechanism}\}
= 1 - \prod_k\left(1 - W_{ik} W_{jk}\right).
\label{eq:eaggl_gene_gene_pair_union_exact}
\end{equation}
When mechanisms are sparse, this union probability is well approximated by the first
order expansion
\begin{equation}
1 - \prod_k\left(1 - W_{ik} W_{jk}\right)
\approx \sum_k W_{ik} W_{jk}.
\label{eq:eaggl_gene_gene_pair_union_approx}
\end{equation}
This approximation motivates a symmetric nonnegative factorization of a pairwise
gene-gene probability matrix.

\paragraph{Shared annotation log-evidence.}
In the PIGEAN annotation model, corrected annotation effects enter the informed gene
prior on the log-odds scale,
\begin{equation}
\log\frac{y_i}{1-y_i}
= \sum_s X_{is}\beta_s + \alpha_i + \log\frac{\pi}{1-\pi},
\label{eq:eaggl_gene_gene_pigean_prior}
\end{equation}
where $X_{is}$ is the value of annotation $s$ for gene $i$, and $\beta_s$ is the
posterior corrected annotation effect. The $\beta_s$ values are therefore measured on
the scale of additive log-odds or log Bayes-factor contributions to gene relevance.
They already include the PIGEAN correction for correlated annotations and the
annotation-specific regularization that accounts for broad or profligate annotations.
The gene-by-gene discovery model therefore does not apply an additional gene-set-size
normalization to $X$ by default.

For a pair of genes $i$ and $j$, EAGGL defines the shared annotation log-evidence
\begin{equation}
L_{ij}
= \sum_s X_{is} X_{js} \beta_s.
\label{eq:eaggl_gene_gene_shared_logbf}
\end{equation}
This quantity is interpreted as the annotation-derived log Bayes factor for the event
that genes $i$ and $j$ are jointly relevant through a common mechanism. Negative
corrected annotation effects are retained in this log-evidence calculation: they lower
pairwise support toward the prior, rather than being truncated before calibration.

\paragraph{Optional profligate-gene correction.}
Some genes have high annotation exposure across many resources and can therefore receive
large shared-annotation scores with many partners. EAGGL leaves the raw pair
log-evidence unchanged by default. When requested, it computes a single exposure feature
for each retained gene from the retained/scored annotation matrix used to construct
$L_{ij}$,
\begin{equation}
a_i = \log\left(1 + \#\{s:X_{is}^{\mathrm{retained}}\ne 0\}\right).
\end{equation}
For each anchor trait, EAGGL fits an ordinary least-squares linear background over
off-diagonal gene pairs,
\begin{equation}
h_{ij} = \alpha + \beta(a_i+a_j).
\end{equation}
The corrected evidence is
\begin{equation}
L_{ij}^{\mathrm{corr}} = L_{ij} - h_{ij}.
\end{equation}
Full read-matrix annotation counts are reported as diagnostics but are not used to fit
the correction.
The downstream conversion to pair probability, excess probability, flooring, diagonal
handling, and anchor aggregation are then identical to the default path. This correction
is opt-in and is intended to remove retained annotation-count exposure, not to change the
definition of the pairwise target otherwise.


\paragraph{Multiple phenotype anchors.}
When gene-by-gene discovery is anchored on multiple phenotypes, EAGGL still estimates
one shared gene-factor basis. For anchor phenotype $t$, let
\begin{equation}
L_{ij}^{(t)}
= \sum_s X_{is} X_{js} \beta_s^{(t)}
\label{eq:eaggl_gene_gene_anchor_logbf}
\end{equation}
denote the anchor-specific shared annotation log-evidence. Each anchor is first
converted exactly as in the single-trait case: the log-evidence is calibrated to a
pairwise same-mechanism probability, transformed to the configured target scale, floored,
and given the configured diagonal policy. Labeling a single-trait input changes only the
reported trait name; it does not change the pair target. Denote the resulting
single-trait target by $T_{ij}^{(t)}$.

The default \texttt{multi} aggregation fits one shared basis against the equal
average of all single-trait targets,
\begin{equation}
\frac{1}{2T}\sum_t \sum_{i,j}
\left(T_{ij}^{(t)} - \sum_k W_{ik}W_{jk}\right)^2
+ \mathrm{ARD}(W).
\label{eq:eaggl_gene_gene_multi_anchor}
\end{equation}
Equivalently, this is symmetric factorization of
$M_{ij}^{\mathrm{multi}} = T^{-1}\sum_t T_{ij}^{(t)}$; anchors with zero target
entries therefore contribute zero values to the average.
The explicit \texttt{any} aggregation uses noisy-OR over the same per-trait targets,
\begin{equation}
M_{ij}^{\mathrm{any}} = 1 - \prod_t \left(1 - T_{ij}^{(t)}\right),
\label{eq:eaggl_gene_gene_any_anchor}
\end{equation}
and then fits the ordinary symmetric objective
$M_{ij}^{\mathrm{any}} \approx \sum_k W_{ik}W_{jk}$. With one anchor trait, both
\texttt{multi} and \texttt{any} reduce exactly to the single-trait target.

\paragraph{Annotation bridge diagnostics.}
For gene-by-gene discovery, EAGGL can optionally decompose the fitted pair-evidence
kernel by annotation after fitting. Each annotation contributes the rank-one kernel
\begin{equation}
K^{(s)} = \beta_s x_s x_s^\top,
\end{equation}
where $x_s$ is the annotation membership vector. Let $\widetilde w_k$ be factor
$k$'s gene-loading column normalized to sum to one. The exact contribution of
annotation $s$ to evidence connecting factors $k$ and $\ell$ is
\begin{equation}
C_{s,k,\ell}
= \widetilde w_k^\top K^{(s)} \widetilde w_\ell
= \beta_s (x_s^\top \widetilde w_k)(x_s^\top \widetilde w_\ell).
\end{equation}
EAGGL reports within-factor mass $\sum_k C_{s,k,k}$, between-factor mass
$\sum_{k\ne\ell} C_{s,k,\ell}$, the corresponding bridge fraction, an effective
number of factors based on normalized annotation-factor overlaps, and a separated
bridge mass that downweights bridges between factors with high weighted-Jaccard gene
overlap. These diagnostics are output-only: they do not alter the learned factors or
the pairwise target.

Suggested annotation exclusions are source-adaptive. EAGGL first marks positive,
active annotations that are broad bridges across low-overlap factors as review
candidates. Automatic exclusion is stricter: an annotation must be a high-impact
diffuse bridge, rank highly both globally and within its source, and come from a
source/anchor group that empirically behaves like a broad bridge source. Broad sources
are identified by a robust source broadness score based on high source bridge burden,
high global bridge burden, and low source quality, rather than by hardcoded library
names. Thus `flag_review` is diagnostic, while `flag_suggest_exclude` is the single
conservative rule used to write the suggested-exclude text file. The same rule is used
before and after exclusion refits; additional suggested exclusions after a refit reflect
changes in the fitted PIGEAN/EAGGL model rather than a separate exclusion mode.

\paragraph{Conversion to pairwise probability.}
Let $\pi_{\mathrm{pair}}$ denote the prior probability that two retained genes are
jointly disease-relevant through the same mechanism, before observing their shared
annotation evidence. The posterior probability is obtained by the usual posterior-odds
conversion,
\begin{equation}
P_{ij}^{\mathrm{same}}
= \sigma\left(\operatorname{logit}(\pi_{\mathrm{pair}}) + L_{ij}\right),
\qquad
\sigma(x)=\frac{1}{1+e^{-x}}.
\label{eq:eaggl_gene_gene_pair_probability}
\end{equation}
Equivalently,
\begin{equation}
\frac{P_{ij}^{\mathrm{same}}}{1-P_{ij}^{\mathrm{same}}}
=
\frac{\pi_{\mathrm{pair}}}{1-\pi_{\mathrm{pair}}}\exp(L_{ij}).
\label{eq:eaggl_gene_gene_pair_odds}
\end{equation}
If annotation effects are stored in a logarithm base other than the natural base, the
implementation first converts $L_{ij}$ to natural-log units before applying
Eq.~\eqref{eq:eaggl_gene_gene_pair_probability}.

The prior $\pi_{\mathrm{pair}}$ controls only the calibration of pairwise probability.
When the user supplies a target mechanism size $S^\star$, a natural default is
\begin{equation}
\pi_{\mathrm{pair}}
\approx \frac{S^\star-1}{N_{\mathrm{active}}-1},
\label{eq:eaggl_gene_gene_pair_prior_from_size}
\end{equation}
where $N_{\mathrm{active}}$ is the number of retained genes. This is the probability
that, conditional on a reference gene, a second randomly selected retained gene is in
the same mechanism under an equal-size, approximately disjoint mechanism model. The
pair prior can also be set directly.

Because a uniform prior background can otherwise be captured by one broad factor, EAGGL
fits the excess pairwise mechanism probability
\begin{equation}
M_{ij}
= \left[\frac{P_{ij}^{\mathrm{same}} - \pi_{\mathrm{pair}}}{1-\pi_{\mathrm{pair}}}\right]_+,
\qquad i\ne j,
\label{eq:eaggl_gene_gene_excess_probability}
\end{equation}
where $[x]_+=\max(x,0)$. Thus $M_{ij}\in[0,1]$ represents the annotation-driven excess
probability that the pair shares a disease-relevant mechanism. The diagonal entries are
excluded or downweighted during fitting, because self-evidence $M_{ii}$ can otherwise
encourage one-gene factors. After this calibration, entries with $M_{ij}<10^{-3}$ are
set to zero by default so that negligible numerical excess probabilities do not create
spurious weak edges in the symmetric factorization target.

\paragraph{Symmetric nonnegative factorization.}
Given the nonnegative pairwise target matrix $M$, EAGGL estimates a nonnegative
gene-by-mechanism loading matrix $W$ by minimizing
\begin{equation}
\mathcal{L}_{\mathrm{sym}}(W,\lambda)
=
\frac{1}{2}
\sum_{i<j} A_{ij}
\left(
M_{ij} - \sum_k W_{ik}W_{jk}
\right)^2
+
\frac{\phi}{2}
\sum_k
\frac{\|W_{\cdot k}\|_2^2 + 2b_0}{\lambda_k}
+
C_{\mathrm{sym}}\sum_k \log\lambda_k
+
\rho\sum_{i,k} W_{ik}.
\label{eq:eaggl_gene_gene_symnmf_objective}
\end{equation}
Here $A_{ij}$ is an optional pairwise fitting weight with $A_{ii}=0$, $\lambda_k$ is
the ARD relevance scale for mechanism $k$, and $\rho$ is an optional within-factor
sparsity penalty. The ARD update is the symmetric analogue of the gene-by-annotation
update,
\begin{equation}
\lambda_k
=
\frac{\frac{1}{2}\|W_{\cdot k}\|_2^2 + b_0}{C_{\mathrm{sym}}},
\qquad
C_{\mathrm{sym}} = \frac{N_{\mathrm{active}}}{2}+a_0+1.
\label{eq:eaggl_gene_gene_lambda_update}
\end{equation}
Nonnegativity is enforced throughout. For probability-scale interpretation, the
implementation can additionally project each row onto the simplex-like constraint
\begin{equation}
0\le W_{ik}\le 1,
\qquad
\sum_k W_{ik}\le 1.
\label{eq:eaggl_gene_gene_row_sum_constraint}
\end{equation}
Under this constraint, Eq.~\eqref{eq:eaggl_gene_gene_pair_union_approx} makes
$\sum_k W_{ik}W_{jk}$ an approximation to the posterior probability that genes $i$ and
$j$ share at least one discovered mechanism beyond the pairwise prior background.

\paragraph{Role of direct gene support.}
Direct genetic support and annotation-derived support are kept separate in this model.
The matrix $M$ is built from the corrected annotation effects $\beta_s$, which already
encode annotation-derived evidence. Multiplying $M$ by the combined gene posterior
probability would double-count the annotation-derived part of the PIGEAN model.
Instead, direct genetic support can be used for reporting, post-fit prioritization, or
optional sensitivity analyses. For example, after fitting $W$, EAGGL can summarize the
direct genetic support of mechanism $k$ by
\begin{equation}
R_k^{\mathrm{direct}}
= \sum_i W_{ik} p_i^{\mathrm{direct}},
\label{eq:eaggl_gene_gene_direct_relevance}
\end{equation}
where $p_i^{\mathrm{direct}}$ is the posterior support derived from direct gene-level
association evidence alone.

\paragraph{Projecting annotations onto gene mechanisms.}
After estimating $W$, all retained annotations are projected onto the learned gene
mechanisms for interpretation. For annotation $s$, EAGGL solves
\begin{equation}
h_s^\star
=
\arg\min_{h_s\ge 0}
\sum_i \eta_i
\left(
X_{is} - \sum_k W_{ik}h_{sk}
\right)^2
+ r(h_s),
\label{eq:eaggl_gene_gene_annotation_projection}
\end{equation}
using the same fixed-basis nonnegative projection machinery used in the rectangular
EAGGL workflow. The fitted annotation loadings $h_{sk}$ provide gene-set labels and
interpretation for each mechanism. Thus, annotations influence discovery through the
pairwise evidence matrix, and additional annotations can be projected afterward to
deepen interpretation without directly redefining the learned gene mechanisms.


\subsection{Phenotype-input anchoring of the gene-by-annotation model}

In some applications, one wishes to condition factorization on several input phenotypes rather than on a single default trait. Let $P_{it}$ denote a nonnegative gene-by-phenotype support matrix and $Q_{jt}$ denote a nonnegative annotation-by-phenotype support matrix derived from phenotype-resolved statistics. EAGGL treats every phenotype with both gene-level and annotation-level support as an anchor. If the complete input phenotype set is $A \subseteq \{1,\ldots,T\}$, then each $t \in A$ defines an anchor user, with
\[
p_i^{(u)} = P_{it},
\qquad
q_j^{(u)} = Q_{jt}.
\]
The weight matrix again takes the additive form of Eq.~\eqref{eq:eaggl_weight_matrix}:
\begin{equation}
\Omega_{ij}^{(A)}
=
\sum_{t \in A}
P_{it} Q_{jt}.
\label{eq:eaggl_pheno_anchor_weights}
\end{equation}
The factorization itself remains the same as Eq.~\eqref{eq:eaggl_basic_factor}. The difference is that the model is now conditioned on a selected phenotype set rather than on a single default trait.

EAGGL no longer exposes separate gene-anchor, gene-set-anchor, or any-anchor public workflows. Those historical modes are retained in git history, but the current public anchoring model is phenotype-input anchoring: users supply one or more phenotype-resolved gene and gene-set evidence tables, EAGGL verifies that the phenotype labels match, and the complete phenotypes jointly define the factorization weights.

\subsection{Implicit factor relevance and any-anchor summaries}

After fitting, EAGGL can compute factor-level relevance scores that summarize how strongly each factor is supported by the selected trait or anchor set. These are downstream summaries derived from the fitted factor basis and are not themselves part of the factorization objective. In the default \texttt{--gene-stats-in} / \texttt{--gene-set-stats-in} workflow, the input statistics define an implicit anchor named \texttt{input\_gene\_stats}; no separate anchor flag is required. The public factor output no longer reports the legacy \texttt{any\_relevance} alias; anchor summaries are exposed through the canonical \texttt{anchor\_any\_joint} and \texttt{anchor\_any\_marginal} columns only when both joint and marginal anchor linkage summaries are available.

For the gene-by-annotation workflow family, let $W \in \mathbb{R}_+^{N \times K}$ denote the fitted gene-by-factor matrix, and let $p^{(u)} \in [0,1]^N$ denote the gene-level support vector for anchor user $u$. The anchor-specific relevance of the factors to user $u$ is defined by bounded nonnegative regression:
\begin{equation}
r^{(u)}
=
\arg\min_{0 \le r \le 1}
\left\|
p^{(u)} - W r
\right\|_2^2.
\label{eq:eaggl_factor_relevance_gene}
\end{equation}

For the gene-set-by-phenotype workflow family, let $Z \in \mathbb{R}_+^{T \times K}$ denote the fitted phenotype-by-factor matrix, and let $p^{(u)} \in [0,1]^T$ denote the phenotype support vector for anchor user $u$. The factor relevance is then defined analogously:
\begin{equation}
r^{(u)}
=
\arg\min_{0 \le r \le 1}
\left\|
p^{(u)} - Z r
\right\|_2^2.
\label{eq:eaggl_factor_relevance_pheno}
\end{equation}

When more than one anchor user is present, EAGGL summarizes the bounded projection relevance of factor $k$ to any selected anchor by a noisy-OR-style aggregation:
\begin{equation}
r_k^{\mathrm{any}}
=
1 -
\prod_{u=1}^U
\left(
1 - r_k^{(u)}
\right).
\label{eq:eaggl_factor_any_relevance}
\end{equation}
This has the intended semantics that $r_k^{\mathrm{any}}$ is a bounded coefficient summary of how strongly factor $k$ links to at least one selected anchor under the relevance model above. In the single-user case, Eq.~\eqref{eq:eaggl_factor_any_relevance} reduces to the anchor-specific relevance itself.

When anchor phenotypes are present and canonical trait linkage is also computed, EAGGL reports \texttt{anchor\_any\_joint} and \texttt{anchor\_any\_marginal} summaries in \texttt{factors.out(.gz)} derived from the canonical linkage calculations in Section~\ref{sec:postfactor_revised}. If only the older joint-style relevance calculation is available, these anchor-any columns are not emitted rather than aliasing marginal to joint. Per-phenotype factor relationships are reported in the canonical long-form \texttt{trait\_factor\_links.out(.gz)} table.

\section{Trait linkage and factor-specific phenotype enrichment}
\label{sec:postfactor_revised}

After fitting the main factorization, EAGGL can relate latent factors to additional human phenotypes in two conceptually distinct ways. The primary public annotation layer is a descriptive \emph{trait linkage} step, which estimates how strongly each phenotype's thresholded high-confidence support profile is represented by each factor. The secondary expert-only layer is an inferential \emph{enrichment} step, which asks whether a phenotype is more concentrated in one factor than expected after accounting for the anchor trait and optional covariates. These two quantities answer different scientific questions and are therefore not expected to coincide exactly.

\subsection{Canonical trait linkage by support-normalized thresholded projection}

The external gene-by-phenotype support table can be very large. In practice, EAGGL therefore consumes a sparse thresholded representation that retains only high-confidence entries. The phenotype profile used for post-factor annotation is therefore not the full unthresholded phenotype surface. It is the thresholded support profile on a single chosen support surface. Publicly, EAGGL defaults to \texttt{--trait-linkage-source combined} and \texttt{--trait-linkage-threshold 1.0}, using a strict threshold rule (\texttt{source value} $> \theta_{\mathrm{trait}}$). Expert overrides can choose \texttt{log\_bf}, \texttt{prior}, or \texttt{auto} fallback resolution.

Let $b$ denote a feature basis shared by the factors and a trait-specific support vector. In the public workflow, $b$ is the gene basis whenever gene-factor loadings are available. EAGGL uses the gene-set basis only as an expert or fallback mode when the gene basis is unavailable or the user explicitly requests \texttt{--project-phenos-from-gene-sets}. Let
\[
W^{(b)} \in \mathbb{R}_+^{F_b \times K}
\]
denote the raw factor-loading matrix on that feature space, where $F_b$ is the number of features in basis $b$, and let
\[
s_t^{(b)} \in \mathbb{R}_+^{F_b}
\]
denote the raw thresholded trait-support profile for phenotype $t$ on that same basis. These raw trait and factor vectors are not interpreted as biological probability distributions and are not required to sum to $1$. Their total support and total mass are meaningful quantities that are preserved separately from the copied normalized vectors used for trait-factor matching.

By default, EAGGL uses a \emph{weighted thresholded} profile: retained source-support values are used for entries that strictly exceed the selected threshold and zero otherwise. As an expert sensitivity analysis, one may instead use a \emph{binary thresholded} profile, in which retained entries are coded as $1$ and absent entries as $0$ before normalization.

For linkage, EAGGL first computes the total thresholded trait support on the full thresholded support surface, before applying the factor mask. Define the total thresholded trait support
\begin{equation}
A_t^{(b)} = \sum_{f=1}^{F_b^{\mathrm{full}}} s_{tf}^{(b,\mathrm{full})},
\label{eq:eaggl_pheno_strength}
\end{equation}
where $s_t^{(b,\mathrm{full})}$ denotes the full thresholded support surface on basis $b$ before masking. Let $\mathcal{M}_b$ denote the retained factor mask on that basis, and let
\[
s_{t,\mathcal{M}}^{(b)}
\]
denote the masked thresholded support profile restricted to the retained factor basis rows. EAGGL then builds a full-space matching target by zeroing support outside $\mathcal{M}_b$, and keeps the normalization denominator as the full thresholded trait strength. Thus the normalized masked trait profile is
\begin{equation}
q_t^{(b)} = \frac{s_{t,\mathcal{M}}^{(b)}}{A_t^{(b)} + \epsilon},
\label{eq:eaggl_pheno_profile_normalized}
\end{equation}
where $\epsilon>0$ is a small constant to avoid division by zero. After masking, $\sum_f q_{tf}^{(b)}$ generally equals the retained support fraction, not $1$:
\[
\sum_f q_{tf}^{(b)} = \frac{\sum_f s_{t,\mathcal{M}f}^{(b)}}{A_t^{(b)} + \epsilon}.
\]
Likewise, each factor basis vector is copied into an internally normalized profile,
\begin{equation}
S_k^{(b)} = \sum_{f=1}^{F_b} W_{fk}^{(b)},
\label{eq:eaggl_factor_total_mass}
\end{equation}
\begin{equation}
b_{\bullet k}^{(b)} = \frac{W_{\bullet k}^{(b)}}{S_k^{(b)} + \epsilon}.
\label{eq:eaggl_factor_basis_normalized}
\end{equation}
The raw factor loadings $W^{(b)}$ and their totals $S_k^{(b)}$ are preserved for reporting; only the copied profiles $b_{\bullet k}^{(b)}$ are used for matching.

EAGGL then reports two linkage summaries from the same internal matching inputs. The \emph{marginal coefficient} of factor $k$ for trait $t$ is the one-factor bounded projection
\begin{equation}
m_{tk}^{(b)}
=
\arg\min_{0 \le c \le 1}
\left\|
q_t^{(b)} - b_{\bullet k}^{(b)} c
\right\|_2^2,
\label{eq:eaggl_trait_linkage_marginal}
\end{equation}
and the marginal overlap companion is the direct shape overlap
\begin{equation}
o_{tk}^{(b)} = \left(q_t^{(b)}\right)^\top b_{\bullet k}^{(b)}.
\label{eq:eaggl_trait_linkage_marginal_overlap}
\end{equation}
and the \emph{joint coefficient} is the all-factor competitive projection
\begin{equation}
j_t^{(b)}
=
\arg\min_{c \ge 0,\ \mathbf{1}^\top c \le 1}
\left\|
q_t^{(b)} - b^{(b)} c
\right\|_2^2.
\label{eq:eaggl_trait_linkage_joint}
\end{equation}
We also report the uncaptured normalized mass
\begin{equation}
\rho_t^{(b)} = \max\!\left(0,\ 1 - \sum_{k=1}^K j_{tk}^{(b)}\right).
\label{eq:eaggl_trait_linkage_residual}
\end{equation}

The coefficient $m_{tk}^{(b)}$ is interpreted as a standalone support-normalized projection coefficient linking phenotype $t$ to factor $k$ when that factor is viewed alone. Because it is divided by the factor self-norm, it remains coefficient-like and can be sensitive to factor breadth. The overlap $o_{tk}^{(b)}$ is a direct shape-overlap metric and is less affected by broad factor profiles. The coefficient $j_{tk}^{(b)}$ is interpreted as the corresponding competitive projection coefficient for factor $k$ after the other factors compete. These coefficients are not exact overlap or captured-support masses. In contrast to a projection normalized only on the retained masked slice, Eqs.~\eqref{eq:eaggl_trait_linkage_marginal} and \eqref{eq:eaggl_trait_linkage_joint} preserve the scale of the full thresholded trait while still solving on the retained factor basis. Broad phenotypes can therefore have large total thresholded support $A_t^{(b)}$ without automatically receiving inflated linkage coefficients when only a small masked subset of their support survives into the factor basis.

For reporting, EAGGL writes:
\begin{itemize}
\item \texttt{trait\_total\_support} = $A_t^{(b)}$
\item \texttt{retained\_trait\_support} = $\sum_f s_{t,\mathcal{M}f}^{(b)}$
\item \texttt{retained\_fraction} = \texttt{retained\_trait\_support} / \texttt{trait\_total\_support}
\item \texttt{trait\_n\_eff} = $\left(\sum_f s_t^{(b)}\right)^2 \big/ \sum_f \left(s_t^{(b)}\right)^2$
\item \texttt{retained\_n\_eff} = $\left(\sum_f s_{t,\mathcal{M}f}^{(b)}\right)^2 \big/ \sum_f \left(s_{t,\mathcal{M}f}^{(b)}\right)^2$
\item \texttt{joint\_coefficient} = $j_{tk}^{(b)}$
\item \texttt{marginal\_coefficient} = $m_{tk}^{(b)}$
\item \texttt{marginal\_overlap} = $o_{tk}^{(b)}$
\item \texttt{joint\_coefficient\_support\_mass} = \texttt{trait\_total\_support} $\times$ \texttt{joint\_coefficient}, a coefficient-scaled support total rather than an exact captured-support mass
\item \texttt{marginal\_coefficient\_support\_mass} = \texttt{trait\_total\_support} $\times$ \texttt{marginal\_coefficient}, a coefficient-scaled support total rather than an exact captured-support mass
\item \texttt{marginal\_overlap\_support\_mass} = \texttt{trait\_total\_support} $\times$ \texttt{marginal\_overlap}, an overlap-scaled support total
\item \texttt{factor\_tier} and \texttt{combined\_mass\_fraction} report the post-fit interpretation tier and global factor mass fraction.
\item \texttt{factor\_total\_mass} = $S_k^{(b)}$ in the factor summary output
\item \texttt{factor\_n\_eff} = $\left(\sum_f W_{fk}^{(b)}\right)^2 \big/ \sum_f \left(W_{fk}^{(b)}\right)^2$
\item \texttt{factor\_top\_share} = $\max_f W_{fk}^{(b)} / S_k^{(b)}$
\item \texttt{factor\_top10\_share} = top-10 loading mass divided by $S_k^{(b)}$
\item \texttt{broad\_factor\_flag} = 1 when \texttt{factor\_n\_eff} $\ge 500$ and \texttt{factor\_top\_share} $\le 0.01$
\end{itemize}

The effective-size diagnostics complement the raw thresholded feature counts. A phenotype can have many retained nonzero genes but still have \texttt{trait\_n\_eff} or \texttt{retained\_n\_eff} close to $1$ when most of its support is concentrated on one or a few genes. These quantities therefore expose concentration of support mass rather than only thresholded breadth.

\paragraph{Interpretation.}
Canonical trait linkage is a descriptive quantity. The marginal coefficients answer:
\begin{quote}
``How much does factor $k$ link to phenotype $t$ when factor $k$ is viewed alone?''
\end{quote}
The joint coefficients answer:
\begin{quote}
``How is the retained high-confidence support profile of phenotype $t$ represented after the factors compete as a mixture of latent mechanisms?''
\end{quote}
These are not p-values, and they should not be interpreted as calibrated posterior probabilities that phenotype $t$ is formally associated with factor $k$.

EAGGL also reports retained-basis diagnostics alongside the linkage scores: total thresholded trait support, retained masked trait support, retained fraction, total thresholded feature count, retained masked feature count, support effective sizes, and a low-retention flag for traits with very small or highly concentrated retained support.

\subsection{Choice of projection basis across workflow families}

For the public gene $\times$ gene-set workflow family, the preferred basis for canonical trait linkage is the gene basis. In that case, $W^{(b)}$ is the gene-by-factor loading matrix and $s_t^{(b)}$ is the thresholded phenotype-specific gene-support vector. If gene-level support is unavailable or gene-level factor loadings are absent, EAGGL falls back to the gene-set basis, in which $W^{(b)}$ is the gene-set-by-factor loading matrix and $s_t^{(b)}$ is the thresholded phenotype-specific gene-set support vector. Thus, the canonical linkage definition is common across workflows: EAGGL always projects one thresholded support profile onto a factor basis defined on one shared feature space, then reports both marginal and joint coefficients from those same internal matching inputs.

\subsection{Factor-specific phenotype enrichment with thresholded binary outcomes}

A second, distinct summary asks whether phenotype $t$ is selectively enriched in genes belonging to factor $k$, after accounting for the main anchor trait and optional covariates. Because the standard phenotype input is thresholded, the default factor-specific enrichment analysis is based on a binary thresholded outcome rather than on truncated continuous support values. Let
\[
h_{it} = \mathbb{I}\{ \text{combined phenotype support for gene } i \text{ and phenotype } t \text{ exceeds } c_0 \}
\]
denote the resulting thresholded phenotype-hit indicator. Let $W_{ik}$ denote the gene-membership weight of factor $k$ for gene $i$, let $a_i^{\mathrm{direct}}$ denote the direct support of the main anchor trait for gene $i$, and let $C_i$ denote optional gene-level covariates. The default model fits, for each phenotype $t$ and factor $k$,
\begin{equation}
h_{it}
=
\alpha_{tk}
+
\beta_{tk} W_{ik}
+
\gamma_t a_i^{\mathrm{direct}}
+
\eta_t^\top C_i
+
\varepsilon_{it}.
\label{eq:eaggl_factor_pheno_regression}
\end{equation}
This regression is fit one factor at a time, i.e. as a marginal anchor-adjusted binary model. The coefficient $\beta_{tk}$ is interpreted as the factor-specific enrichment of phenotype $t$ beyond what is already captured by the anchor trait's direct support.

\paragraph{Interpretation.}
Unlike the canonical linkage weights in Eqs.~\eqref{eq:eaggl_trait_linkage_marginal} and \eqref{eq:eaggl_trait_linkage_joint}, the coefficients in Eq.~\eqref{eq:eaggl_factor_pheno_regression} are incremental effects. They answer the question:
\begin{quote}
``Is phenotype $t$ enriched in factor $k$ beyond the anchor trait's direct gene support?''
\end{quote}
This is an inferential or hypothesis-testing quantity rather than a descriptive loading.

\paragraph{Recommended default and scope.}
The binary thresholded model is the default because it is defensible under sparse thresholded phenotype input and easier to interpret than a joint all-factor fit. EAGGL reports robust uncertainty by default, typically using HC3-type standard errors or chromosome-clustered uncertainty when location information is available. The resulting p-values are still approximate, because the thresholded phenotype-hit process is derived from a combined-support quantity that already reflects gene-set-mediated indirect support.

An expert may instead fit: (i) a marginal binary model without an anchor covariate; (ii) a joint all-factor binary model with direct anchor adjustment; or (iii) legacy continuous direct/combined regressions. Using combined anchor support as the covariate is also an expert-only approximation, because then both the outcome and the covariate inherit related gene-set-mediated structure.

When multiple factor-specific enrichment models are requested in one run, EAGGL appends them into a single output table and tags each row with the model name, factor-model scope, outcome surface, and anchor-adjustment setting so that marginal, joint, and legacy analyses are not conflated downstream.

\subsection{Relationship between canonical trait linkage and factor-specific enrichment}

The two phenotype summaries serve complementary roles:
\begin{itemize}
\item \textbf{Canonical trait linkage} gives bounded marginal and joint support-normalized projection coefficients describing how strongly the retained high-confidence phenotype profile links to each factor. This is the primary annotation layer for interpreting factors with external phenotypes.
\item \textbf{Factor-specific enrichment regression} tests whether that phenotype profile is unusually concentrated in one factor beyond what is already captured by the anchor trait and optional covariates. This is a secondary, expert-oriented summary.
\end{itemize}

Consequently, the main phenotype labels and top-phenotype annotations for a factor should be based on canonical trait-linkage weights, whereas regression coefficients and associated significance measures should be interpreted as downstream enrichment diagnostics rather than as the primary phenotype annotation of the factor.

\subsection{Repeated restarts, consensus factorization, and structural selection of $\phi$}

The EAGGL objective is non-convex, and therefore different random initializations can converge to different local optima. For this reason, EAGGL is typically fit under repeated random restarts. The simplest use of repeated restarts is to retain the fitted solution with the best objective value. A more stringent alternative is consensus factorization, in which factors are first matched across runs and then aggregated only if they are reproduced consistently across restarts. Let $H^{(r)}$ denote the annotation loading matrix from restart $r$ after ARD thresholding. In consensus mode, EAGGL normalizes the factor columns, aligns factors across restarts by maximum total similarity, and aggregates matched factors only when they appear in a sufficient fraction of runs.

The hyperparameter $\phi$ also affects the effective number and sparsity of retained factors, even though ARD is responsible for pruning irrelevant components. EAGGL therefore treats $\phi$ as a structural model-selection parameter. The default automatic selection objective is a composite score over interpretable candidate-level metrics. For candidate $c$ and component $j$, let $q_{cj} \in [0,1]$ be the component score and let $a_{cj}$ indicate that the component is available. EAGGL selects by
\begin{equation}
Q_c = \frac{\sum_j a_{cj}\omega_j q_{cj}}{\sum_j a_{cj}\omega_j},
\label{eq:eaggl_phi_composite_score}
\end{equation}
where unavailable components are skipped rather than treated as zero. The default weights are 0.15 for factor size, 0.15 for non-overlap, 0.15 for entity concentration, 0.25 for high-priority coverage, 0.10 for reconstruction, 0.10 for within-factor coherence, 0.05 for factor balance, and 0.05 for annotation bridge quality control. Ties within the configured tolerance prefer fewer factors, then higher coverage, then lower $\phi$.

The factor-size component compares each factor's capped nonnegative gene loading mass $S_k=\sum_i \min(W_{ik},c)$ with a target $S_0$:
\begin{equation}
q^{\mathrm{size}}_k = \left[1 + \left(\frac{\log_2((S_k+\epsilon)/(S_0+\epsilon))}{w}\right)^2\right]^{-1},
\end{equation}
and averages this score over factors. Non-overlap uses the 90th percentile of soft weighted-Jaccard overlaps between factor loading vectors. Entity concentration scores whether genes and, when available, annotations concentrate on one factor rather than bridging many factors. Coverage scores whether high-priority genes and annotations, ranked by the resolved PIGEAN evidence columns, have at least one loading above the configured coverage threshold. Reconstruction and coherence use the actual discovery target: $WW^T$ for gene-by-gene discovery and $HW^T$ for gene-by-annotation discovery, respecting rectangular anchor weights when available. Factor balance is normalized entropy of factor gene mass. Annotation bridge QC is included only when annotation bridge metrics are available.

A legacy selector remains available. In legacy mode, users specify a target median effective gene support for primary factors, and EAGGL searches for a $\phi$ whose fitted mechanisms have that approximate gene-module size. Within each modal restart, EAGGL measures redundancy on the gene loading matrix when available, with fallback to the gene-set and phenotype loading matrices. These phi-selection metrics are computed over the primary factor set by default so that low-mass ARD-tail columns do not dominate model selection; an audit option can instead score all fitted columns. Candidate values are explored around the initial $\phi$ by target-aware bracketing. The selected $\phi$, candidate diagnostics, component scores, weights, weighted contributions, search thresholds, factor-mass summaries, and convergence diagnostics are written to the run log and parameter output. Per-factor diagnostics for each tested $\phi$ can also be written for audit and review.

EAGGL separates fitted factors from reported factors. Let $m_k$ denote the combined nonnegative gene and annotation loading mass for factor $k$, and let $\pi_k=m_k/\sum_{\ell}m_{\ell}$. A factor is classified as primary when $\pi_k \ge 0.005$, secondary when $0.001 \le \pi_k < 0.005$, and filtered otherwise. The default user-facing factor and cluster outputs print only primary factors, because low-mass secondary and filtered factors usually represent the ARD tail rather than interpretable mechanisms. This reporting scope is controlled separately from the phi-selection metric scope. The exhaustive factor-metrics outputs remain available for auditing all fitted columns. The \texttt{factors.out(.gz)} table records both \texttt{factor\_tier} and \texttt{combined\_mass\_fraction} next to \texttt{lambda}; widening the output scope can include secondary or all factors without changing the fitted model or renumbering factor identifiers.

\subsection{Optimizations for scaling EAGGL to large phenotype panels}

Although the model is defined in terms of weighted nonnegative factorization, a direct implementation that explicitly constructs separate dense matrices for every anchor user would be inefficient. EAGGL therefore uses the weighted objective of Eq.~\eqref{eq:eaggl_weighted_likelihood} directly, rather than materializing large user-specific tensors. This preserves the model while substantially reducing memory use.

Similarly, phenotype panels and annotation collections can be large enough that a naive factorization would be dominated by nearly null rows or by many near-duplicate annotations. EAGGL therefore distinguishes three annotation sets during factor fitting: retained annotations, discovery annotations, and projected annotations. Retained annotations are the rows that survive ordinary hard support and quality filters. Discovery annotations are a redundancy-balanced subset of family leaders used to learn the shared factor basis $W$. Projected annotations are the full retained set, obtained afterward by fixed-$W$ nonnegative projection.

In the default discovery mode, retained annotations are processed in support-ranked order and greedily assigned to similarity-defined families. Let $F_c$ denote discovery family $c$, let $r_c$ denote its leader, and let $s(j,r_c) \in [0,1]$ be the same clipped annotation similarity used for assignment. A candidate annotation starts a new family when its maximum similarity to existing leaders is less than the discovery similarity threshold $\rho_{\mathrm{disc}}$; otherwise it is assigned to the nearest leader. For singleton families $m_{c,\mathrm{eff}}=1$. For larger families,
\begin{equation}
\bar s_c = \frac{1}{|F_c|-1}\sum_{j \in F_c \setminus \{r_c\}} s(j,r_c),
\qquad
m_{c,\mathrm{eff}} = \frac{|F_c|}{1 + (|F_c|-1)\bar s_c}.
\label{eq:eaggl_discovery_effective_size}
\end{equation}
The default effective-size discovery support for anchor user $u$ is
\begin{equation}
\tilde q_c^{(u)} = q_{r_c}^{(u)} m_{c,\mathrm{eff}}.
\label{eq:eaggl_discovery_effective_weight}
\end{equation}
Thus exact duplicates contribute approximately one representative unit, whereas weakly redundant families can contribute closer to their family size. The conservative \texttt{log\_effective\_size} mode uses $q_{r_c}^{(u)}(1+\log m_{c,\mathrm{eff}})$, and \texttt{none} uses the representative support $q_{r_c}^{(u)}$ without a family-size multiplier. EAGGL records \texttt{in\_discovery}, \texttt{discovery\_family\_size}, \texttt{discovery\_family\_mean\_similarity}, \texttt{discovery\_family\_effective\_size}, and \texttt{discovery\_weight} in cluster outputs so the discovery objective can be audited directly.

Because discovery support weights change the data term scale, the implementation computes internal ARD scaling quantities from the row-weighted discovery matrix rather than the raw discovery matrix. Equivalently, if $V$ is the discovery matrix and $D$ is the diagonal matrix of discovery row weights, scale-dependent quantities use $D^{1/2}V$. This keeps the effective $\phi$ scaling, ARD prior floor, and related heuristics aligned with the weighted objective. Finally, output cluster tables first apply the requested factor output scope and then omit gene and gene-set rows whose maximum raw factor loading among the printed factors is below the configurable \texttt{--cluster-row-min-max-loading} threshold; these are reporting filters and do not alter the fitted factor matrices.

\subsection{Interpretation and labeling of latent mechanisms}

The outputs of EAGGL are the fitted factor loadings over genes, annotations, and, when applicable, phenotypes. These loadings constitute the statistical model output. Human-readable labels for factors are downstream summaries rather than part of the factorization model itself.

If $\mathcal{S}_k$, $\mathcal{G}_k$, and $\mathcal{P}_k$ denote the highest-weight annotations, genes, and phenotypes for factor $k$, then a label may be viewed abstractly as a post hoc summary function
\[
L_k = \mathcal{L}(\mathcal{S}_k, \mathcal{G}_k, \mathcal{P}_k).
\]
Such labels may be generated deterministically from the top terms or by a separate language-model-based summarization step. In either case, the labeling procedure is downstream of the core statistical model and should be interpreted as an aid to biological interpretation rather than as part of the inference procedure itself. When users provide a designated labeling library, EAGGL restricts the annotation candidates in $\mathcal{S}_k$ to that library before ranking terms for the label; if a factor has no positive-loading annotation in the designated library, the label falls back to the full annotation set so that every factor remains labeled. This filtering changes only the post hoc label summary and not the fitted factors or output loadings.

\end{document}

\subsection{Simplified trait-factor linkage}

For downstream trait interpretation, EAGGL reports two factor-to-trait summaries. First, each factor is treated as a weighted gene set after setting small gene loadings below a configurable threshold to zero. The resulting factor gene sets are scored against trait support surfaces, yielding factor-level \texttt{beta}, \texttt{beta\_uncorrected}, \texttt{beta\_tilde}, standard error, $z$, and $p$ columns in \texttt{trait\_factor\_links.out}. Second, EAGGL projects each trait support vector onto the factor gene-loading basis by nonnegative least squares and reports the coefficient as \texttt{nnls\_loading}. These two quantities replace the older capture and lift columns in the default trait-factor linkage output.

Cluster output tables contain the raw factor loading for each row and a cosine column for each factor. For row loading vector $v_i$, the reported \texttt{Cosine\_Factor}$k$ value is $v_{ik}/\|v_i\|_2$, i.e. the cosine similarity between the row's full factor-loading vector and the one-hot indicator for factor $k$.
